% !TeX root = ../main.tex
% !TeX encoding = UTF-8
\subsection{Ablation Study}\label{sec:exp-ablation}

We evaluated the three residual pathways with Chengdu ablations. A1 removes the self-evolution expert, A2 removes the directed time-lagged expert, and A3 removes BCMS. Each variant was retrained for up to 200 epochs with an early-stopping patience of 20 epochs. In A1 and A2, the remaining expert weights were renormalized after setting the removed weight to zero. Testing used the positive-target mask and AVG12 aggregation defined in Section~\ref{sec:exp-baselines-metrics}. Joint MAE and RMSE are the equal-weight averages of the inflow and outflow AVG12 values.

% !TeX root = ../main.tex
\begin{table}[!htbp]
  \centering
  \caption{Chengdu ablation results (AVG12)}
  \label{tab:serr-expert-ablation}
  \begingroup
  \linespread{1}\fontsize{9}{11}\selectfont
  \setlength{\tabcolsep}{3pt}
  \renewcommand{\arraystretch}{1.08}
  \begin{tabular*}{\linewidth}{@{\extracolsep{\fill}}lrrrrrr@{}}
    \toprule
    Model configuration & \shortstack{Inflow\\MAE} & \shortstack{Inflow\\RMSE} & \shortstack{Outflow\\MAE} & \shortstack{Outflow\\RMSE} & \shortstack{Joint\\MAE} & \shortstack{Joint\\RMSE} \\
    \midrule
    RBSRN & \ablationpolished{\textbf{0.7084}} & \ablationpolished{\textbf{1.2935}} & \ablationpolished{\textbf{0.7087}} & \ablationpolished{\textbf{1.2235}} & \ablationpolished{\textbf{0.7085}} & \ablationpolished{\textbf{1.2585}} \\
    w/o Self-evolution expert (A1) & \ablationpolished{0.7129} & \ablationpolished{1.3037} & \ablationpolished{0.7121} & \ablationpolished{1.2270} & \ablationpolished{0.7125} & \ablationpolished{1.2653} \\
    w/o Directed time-lagged expert (A2) & \ablationpolished{0.7105} & \ablationpolished{1.3048} & \ablationpolished{0.7099} & \ablationpolished{1.2264} & \ablationpolished{0.7102} & \ablationpolished{1.2656} \\
    w/o BCMS (A3) & \ablationpolished{0.7109} & \ablationpolished{1.3014} & \ablationpolished{0.7102} & \ablationpolished{1.2276} & \ablationpolished{0.7105} & \ablationpolished{1.2645} \\
    \bottomrule
  \end{tabular*}
  \endgroup
\end{table}

\FloatBarrier

Table~\ref{tab:serr-expert-ablation} shows that the full RBSRN had the lowest error on all six metrics. Removing any branch increased MAE and RMSE for both flows. Removing the self-evolution expert increased joint MAE by 0.56\%, the largest increase among the variants. Removing the directed time-lagged expert increased joint RMSE by 0.56\%, compared with 0.54\% for removing self-evolution. The two experts therefore contribute through different error measures, with self-evolution having the larger effect on joint MAE.

Removing BCMS increased joint MAE and RMSE by 0.28\% and 0.48\%, respectively, with higher errors for both flows. The full model performed best when self-evolution, directed time-lagged association, and background-conditioned message selection were combined.
\FloatBarrier
